\documentclass[conference]{IEEEtran}
\usepackage[utf8]{inputenc}
\usepackage[spanish,es-tabla]{babel}
\usepackage{amsmath,amssymb,amsfonts}
\usepackage{algorithmic}
\usepackage{graphicx}
\usepackage{textcomp}
\usepackage{xcolor}
\usepackage{cite}
\usepackage{array}
\usepackage{xurl}
\usepackage[hidelinks]{hyperref}
\usepackage[section]{placeins}
\raggedbottom

\title{\textbf{Optimización de la gestión de recursos críticos hospitalarios: integración de pronósticos epidemiológicos y programación lineal ante brotes estacionales}\\[0.5em]
\large Trabajo Final Integrador}

\author{
    \IEEEauthorblockN{Bellor Maria E., Diaz Dylan S., Jimenez Fernando E., Veliz Ignacio M., Zato Sosa Celina A.}
    \IEEEauthorblockA{Universidad Tecnológica Nacional, Facultad Regional Tucumán\\
    Investigación Operativa\\
    \textbf{Profesores:} Dra. Gramajo Guadalupe, Ing. Rojas Cristina, Ing. Fraga Alvaro\\
    Tucumán, Julio de 2026}
}

\begin{document}

\maketitle

\begin{abstract}
    Este trabajo propone un sistema de soporte de decisiones fundamentado en Investigación Operativa para optimizar la gestión logística hospitalaria en Tucumán ante la saturación provocada por brotes estacionales de dengue, bronquiolitis e influenza. La metodología integra un modelo predictivo Holt-Winters Multiplicativo para proyectar la demanda semanal y un modelo de Programación por Metas Lexicográfica estructurado en subetapas secuenciales que segmentan la atención según triaje clínico (cuidados intensivos, hospitalización general y atención ambulatoria). El algoritmo prioriza sistemáticamente la cobertura médica de los casos de mayor gravedad, garantizando la atención integral de los pacientes críticos e internados, al tiempo que la etapa financiera minimiza el sobrecosto operativo sin degradar los logros médicos alcanzados. Los resultados demuestran que la saturación asistencial no se origina en la infraestructura de camas ni en el personal de salud, sino en cuellos de botella logísticos de insumos descartables específicos, ofreciendo una herramienta transparente para la planificación sanitaria proactiva.
\end{abstract}

\begin{IEEEkeywords}
    Investigación Operativa, Programación por Metas Lexicográfica, Pronósticos Epidemiológicos, Holt-Winters, Asignación de Recursos Hospitalarios.
\end{IEEEkeywords}

\section{Fundamentación}
Los sistemas de salud pública en Tucumán enfrentan tensiones extremas ante la coexistencia estacional de dengue, bronquiolitis e influenza, provocando la frecuente saturación hospitalaria y el desabastecimiento de insumos críticos \cite{ref1}. Esta volatilidad epidemiológica exige pasar de una respuesta reactiva a una planificación sanitaria proactiva.

Para evitar el colapso, es vital anticipar la demanda de pacientes. El análisis de series temporales mediante la Suavización Exponencial Triple de Holt-Winters (variante Multiplicativa) ha demostrado ser altamente efectivo en la predicción de brotes infecciosos. Su elección se justifica en que la amplitud de los contagios masivos estacionales de dengue, bronquiolitis e influenza crece proporcionalmente al nivel general de la epidemia, permitiendo capturar de manera precisa las tendencias y los marcados patrones estacionales \cite{ref2, ref3}.

Sin embargo, este pronóstico debe acompañarse de una gestión logística eficiente. La aplicación de la Investigación Operativa en sistemas de salud, mediante modelos de Programación por Metas (Goal Programming), permite optimizar la asignación de recursos escasos, mejorando la utilización de la infraestructura y reduciendo costos operativos al anticiparse a la saturación \cite{ref4, ref5}.

Por lo tanto, este trabajo se fundamenta en la sinergia de ambas herramientas. Los pronósticos de Holt-Winters estimarán la demanda futura, y un modelo de Programación por Metas Lexicográfica optimizará la asignación semanal de recursos hospitalarios finitos, priorizando la atención de los casos de mayor gravedad clínica de forma innegociable frente al costo, para garantizar equidad, eficiencia y minimizar la mortalidad.

\section{Objetivos}
Ante la histórica ineficiencia distributiva del sistema sanitario y la consecuente saturación hospitalaria durante picos epidemiológicos, resulta imperativo adoptar herramientas de planificación proactiva que garanticen una atención equitativa y oportuna \cite{ref6}. En respuesta a esta necesidad, se define:

\subsection{Objetivo Principal}
Desarrollar un sistema de soporte de decisiones fundamentado en modelos de Investigación Operativa para optimizar la gestión logística y la asignación de recursos críticos en hospitales públicos de la provincia de Tucumán frente a la demanda estacional aguda por brotes de dengue, bronquiolitis e influenza.

\subsection{Objetivos Específicos}
\begin{itemize}
    \item Proyectar la demanda semanal de atención médica para las patologías en estudio mediante el análisis predictivo de series temporales, implementando el método de Suavización Exponencial Triple (Holt-Winters Multiplicativo) para capturar tendencias y estacionalidades.
    \item Formular un modelo de Programación por Metas Lexicográfica de carácter estático que sistematice la asignación de capacidades limitadas, incluyendo camas de internación, unidades de cuidados críticos (UTI), terapias de oxígeno, kits médicos y horas de personal.
    \item Minimizar los costos operativos a través de la optimización del algoritmo matemático, garantizando la eficiencia financiera sin comprometer la cobertura asistencial.
    \item Maximizar la cantidad de pacientes atendidos integrando un factor de ponderación clínica que priorice sistemáticamente los casos de mayor letalidad, garantizando el cumplimiento de los protocolos de resguardo ético-médico.
\end{itemize}

\section{Marco Conceptual}
\begin{itemize}
    \item \textbf{Patología epidémica estacional:} Enfermedad infecciosa aguda con incrementos masivos y cíclicos anuales (dengue, bronquiolitis e influenza).
    \item \textbf{Demanda epidemiológica aguda:} Picos repentinos de pacientes que amenazan con saturar la capacidad instalada y agotar los presupuestos.
    \item \textbf{Recurso crítico asistencial:} Elemento finito e indispensable (camas de internación, UTI, oxigenoterapia, kits médicos y personal de salud).
    \item \textbf{Capacidad neta hospitalaria:} Disponibilidad operativa real, sujeta a restricciones estrictas de presupuesto, espacio y protocolos.
    \item \textbf{Planificación sanitaria proactiva:} Estrategia para anticipar necesidades futuras y evitar el colapso hospitalario frente a crisis epidemiológicas.
    \item \textbf{Ponderación clínica (Triage):} Priorización matemática de atención basada en la severidad y la letalidad histórica de cada paciente.

    \item \textbf{Logística de Contingencia Sanitaria:} Abastecimiento y distribución eficiente de recursos para sostener la operatividad ininterrumpida.
    \item \textbf{Costo Operativo Hospitalario:} Erogación financiera base y extraordinaria requerida para sostener la infraestructura y el capital humano.
\end{itemize}

\section{Marco Teórico}
\begin{itemize}
    \item \textbf{Pronóstico:} Técnica analítica de la Investigación Operativa que utiliza datos históricos y modelos matemáticos para proyectar escenarios futuros, reduciendo la incertidumbre del sistema \cite{ref7, ref11}.
    \item \textbf{Suavización Exponencial Triple (Holt-Winters):} Método que descompone y proyecta simultáneamente nivel, tendencia y estacionalidad. Es indispensable para estructurar alertas tempranas en enfermedades estacionales \cite{ref9}.
    \item \textbf{Investigación Operativa (IO):} Disciplina cuantitativa que optimiza el rendimiento de un sistema sujeto a restricciones, proveyendo el soporte analítico para readecuar recursos hospitalarios ante brotes \cite{ref10}.
    \item \textbf{Programación por Metas (Goal Programming):} Técnica matemática de optimización que busca alcanzar múltiples metas simultáneas minimizando sus desvíos \cite{ref8}.
    \item \textbf{Prioridad Lexicográfica:} Enfoque de optimización multiobjetivo que establece un orden de prioridades absolutas entre metas en conflicto. Las metas de mayor jerarquía (ética médica) se cumplen primero, sin que un mejor desempeño en metas inferiores (costos) pueda empeorar lo ya logrado \cite{ref11, refGP}.
\end{itemize}


\section{Marco Contextual}
\subsection{Alcance Metodológico y Carácter General del Modelo}
El sistema propuesto constituye un marco matemático genérico y adaptable de soporte a las decisiones logístico-asistenciales, concebido como una herramienta de planificación proactiva para prever el uso más eficiente de los recursos hospitalarios ante crisis epidemiológicas. Se declara explícitamente que el modelo de optimización es de carácter estático y monoperíodo; evalúa un horizonte de decisión semanal aislado sin arrastrar internaciones de períodos previos, inventarios acumulados ni órdenes de compra futuras.

Es fundamental destacar que este modelo no constituye una regla matemática rígida de cumplimiento obligatorio ni una directiva automatizada de asignación taxativa. Funciona como un instrumento cuantitativo de orientación estratégica cuya aplicación práctica estará siempre supeditada al criterio profesional médico, al juicio clínico individualizado y a los principios éticos y morales que rigen a la sociedad. La arquitectura del modelo es modular y secuencial: se alimenta en primera instancia por los vectores de demanda esperada ($D_i$) provistos por el modelo predictivo determinístico de pronósticos (Holt-Winters), y posteriormente por la matriz de parámetros operativos ($K_{\text{max},r}, B_r, T_{i,r}, C_r$) específicos de cualquier efector de salud que requiera desplegarlo.

Para validar el funcionamiento práctico y la robustez del algoritmo ante un escenario epidemiológico agudo real, se utilizó como prototipo de estudio de referencia al \textbf{Hospital Centro de Salud Zenón J. Santillán} \cite{refMSPT2}, efector de tercer nivel de la provincia de Tucumán. La parametrización del caso de prueba integra una combinación de datos institucionales oficiales confirmados (capacidad física de camas, normativas salariales y personal de plantilla) y valores estimados basados en directrices técnicas de la OMS/OPS (matriz de consumo tecnológico $T_{i,r}$ y costos de insumos de contingencia). Los supuestos, fórmulas y desgloses de dichos parámetros estimados se detallan explícitamente en las planillas de soporte del trabajo \cite{refGP}.

\subsection{Contexto Demográfico y Población de Riesgo}
El modelo se enmarca en la provincia de Tucumán. Según datos del Censo Nacional 2022 \cite{ref12}, de sus 1.703.186 habitantes, el 42\% depende de manera exclusiva del subsistema público de salud (SIPROSA). La segmentación etaria es fundamental para estimar la demanda potencial asimétrica ante el brote:
\begin{table}[htbp]
    \caption{Distribución demográfica de riesgo y dependencia del sistema público en Tucumán (basado en \cite{ref12})}
    \label{tab:demografia}
    \centering
    \setlength{\tabcolsep}{3pt}
    \begin{tabular}{|>{\raggedright\arraybackslash}m{1.8cm}|>{\raggedright\arraybackslash}m{1.8cm}|>{\raggedright\arraybackslash}m{1.4cm}|>{\raggedright\arraybackslash}m{2.1cm}|}
        \hline
        \multicolumn{1}{|>{\centering\arraybackslash}m{1.8cm}|}{\textbf{Grupo Poblacional}} & \multicolumn{1}{>{\centering\arraybackslash}m{1.8cm}|}{\textbf{Patología de Mayor Riesgo}} & \multicolumn{1}{>{\centering\arraybackslash}m{1.4cm}|}{\textbf{Población Total (Tucumán)}} & \multicolumn{1}{>{\centering\arraybackslash}m{2.1cm}|}{\textbf{Demanda Potencial (Dep. Pública $\approx$42\%)}} \\
        \hline
        Primera Infancia (0 a 4 años)                                                       & Bronquiolitis (VSR)                                                                        & $\approx$ 142.500 hab.                                                                     & $\approx$ 59.850 lactantes y niños                                                                              \\
        \hline
        Población General (15 a 64 años)                                                    & Dengue                                                                                     & $\approx$ 1.085.000 hab.                                                                   & $\approx$ 455.700 jóvenes y adultos                                                                             \\
        \hline
        Adultos Mayores (65 años y más)                                                     & Influenza A / Neumonía                                                                     & $\approx$ 186.200 hab.                                                                     & $\approx$ 78.200 adultos mayores                                                                                \\
        \hline
    \end{tabular}
\end{table}

\subsection{Escenario de Máxima Tensión Epidemiológica (Referencia Histórica)}
Para parametrizar el comportamiento de la demanda ($D_i$) en crisis, se toma el escenario histórico del ciclo 2023-2024 en Tucumán, donde confluyeron el mayor brote histórico de dengue y un aumento temprano de Infecciones Respiratorias Agudas Bajas (IRAB). La saturación máxima se registró entre las Semanas Epidemiológicas (SE) 13 y 17 del 2024:
\begin{table}[htbp]
    \caption{Registro histórico del pico de crisis multi-epidémica en Tucumán (valores de referencia basados en \cite{ref13})}
    \label{tab:pico_critico}
    \centering
    \setlength{\tabcolsep}{3pt}
    \begin{tabular}{|>{\raggedright\arraybackslash}m{1.7cm}|>{\raggedright\arraybackslash}m{1.3cm}|>{\raggedright\arraybackslash}m{2.0cm}|>{\raggedright\arraybackslash}m{2.0cm}|}
        \hline
        \multicolumn{1}{|>{\centering\arraybackslash}m{1.7cm}|}{\textbf{SE (Pico Crítico)}} & \multicolumn{1}{>{\centering\arraybackslash}m{1.3cm}|}{\textbf{Casos Dengue}} & \multicolumn{1}{>{\centering\arraybackslash}m{2.0cm}|}{\textbf{Bronquiolitis ($<$2 años)}} & \multicolumn{1}{>{\centering\arraybackslash}m{2.0cm}|}{\textbf{ETI / Influenza}} \\
        \hline
        SE 13 (Fin Marzo)                                                                   & 7.290                                                                         & 64                                                                                         & 354                                                                              \\
        \hline
        SE 14 (Inic. Abril)                                                                 & 8.333                                                                         & 69                                                                                         & 375                                                                              \\
        \hline
        SE 15 (Med. Abril)                                                                  & 8.160                                                                         & 72                                                                                         & 328                                                                              \\
        \hline
        SE 16 (Fin Abril)                                                                   & 6.878                                                                         & 117                                                                                        & 437                                                                              \\
        \hline
        SE 17 (Inic. Mayo)                                                                  & 6.708                                                                         & 114                                                                                        & 531                                                                              \\
        \hline
    \end{tabular}
\end{table}

Para alimentar y calibrar el modelo predictivo en este contexto de crisis, se definieron criterios metodológicos específicos para la delimitación de las ventanas de datos históricos de cada patología:
\begin{itemize}
    \item \textbf{Dengue (Ciclos 2023--2024):} Se limitó a este periodo debido a que representó el evento de máxima crisis sanitaria en la historia de Tucumán. Incorporar años previos de baja circulación endémica distorsionaría la estacionalidad y restaría capacidad de alerta temprana ante brotes agudos y de comportamiento vectorial altamente dependiente de cepas e indicadores climáticos.
    \item \textbf{Bronquiolitis e Influenza (Ciclos 2023--2025):} Se empleó una ventana de tres años. Estas patologías presentan un comportamiento endoepidémico invernal regular y cíclico, por lo que una muestra más amplia asegura la estabilidad estadística necesaria para modelar las Infecciones Respiratorias Agudas Bajas (IRAB).
\end{itemize}

\subsection{Parámetros Clínicos: Hospitalización, Severidad y Letalidad}
La conversión de casos en requerimientos de insumos emplea tasas de consumo tecnológico ($T_{i,r}$) y ponderadores clínicos de gravedad ($\mu_i$), estructurados bajo directrices de la Organización Mundial de la Salud (OMS) y la Organización Panamericana de la Salud (OPS). El modelo filtra la demanda bruta según tasas de hospitalización esperadas e índices de letalidad:

Es preciso aclarar que el coeficiente $\mu_i$ constituye una escala ordinal de severidad que expresa el orden de prioridad clínica entre patologías dentro de una misma modalidad asistencial (baja/moderada = 1 para Dengue, alta = 2 para Bronquiolitis y muy alta = 3 para Influenza). Esta asignación numérica representa un orden jerárquico relativo y no una relación de proporcionalidad cardinal; en particular, no significa que la influenza sea exactamente tres veces más grave que el dengue ni que la bronquiolitis sea el doble, sino que proporciona al algoritmo un criterio de jerarquización ordinal para dar mayor importancia a los déficits asociados a cuadros de mayor severidad dentro de cada modalidad.

\begin{table}[htbp]
    \caption{Indicadores clínicos de severidad y requerimientos asistenciales (basado en \cite{ref14, ref15})}
    \label{tab:indicadores_clinicos}
    \centering
    \setlength{\tabcolsep}{3pt}
    \begin{tabular}{|>{\raggedright\arraybackslash}m{1.5cm}|>{\raggedright\arraybackslash}m{1.8cm}|>{\raggedright\arraybackslash}m{1.8cm}|>{\raggedright\arraybackslash}m{2.1cm}|}
        \hline
        \multicolumn{1}{|>{\centering\arraybackslash}m{1.5cm}|}{\textbf{Patología ($i$)}} & \multicolumn{1}{>{\centering\arraybackslash}m{1.8cm}|}{\textbf{Hosp. (Camas)}} & \multicolumn{1}{>{\centering\arraybackslash}m{1.8cm}|}{\textbf{Cuidados Críticos (UTI/ARM)}} & \multicolumn{1}{>{\centering\arraybackslash}m{2.1cm}|}{\textbf{Prioridad / Letalidad ($\mu_i$)}} \\
        \hline
        1. Dengue                                                                         & 5\% - 7\% (Crítica / Alarma)                                                   & $<$ 1\% (Dengue Grave)                                                                       & Bajo / Mod. ($\mu_1 = 1$)                                                                        \\
        \hline
        2. Bronquiolitis (VSR)                                                            & 10\% - 15\% (Hipoxemia)                                                        & 3\% - 5\% (UTI)                                                                              & Alto ($\mu_2 = 2$)                                                                               \\
        \hline
        3. Influenza A / ETI                                                              & 8\% - 12\% (Riesgo)                                                            & 4\% - 6\% (Insuf. Resp.)                                                                     & Muy Alto ($\mu_3 = 3$)                                                                           \\
        \hline
    \end{tabular}
\end{table}

\subsection{Operativización Financiera y Restricciones de Capacidad}
Los efectores de tercer nivel del SIPROSA operan con capacidad instalada rígida. El modelo calcula la capacidad neta deduciendo la carga base histórica ocupada por patologías no epidémicas (entre 60\% y 70\% de camas). Por su parte, los costos unitarios ($C_r$) y el presupuesto extraordinario semanal ($PresupuestoMax$) se ingresan dinámicamente según valores vigentes de mercado y asignaciones ministeriales, garantizando la aplicabilidad temporal del modelo al calcular la frontera de eficiencia óptima.

\subsection{Origen de Datos y Metodología de Digitalización}
Los datos cuantitativos semanales reales ($y_t$) se obtuvieron de los reportes oficiales de la sala de situación epidemiológica del SIPROSA y de los Boletines Epidemiológicos Nacionales, cubriendo desde marzo de 2023 hasta enero de 2026.

Debido a que múltiples reportes históricos solo presentaban las curvas epidemiológicas en formato gráfico no editable (PDFs), se utilizó la herramienta digitalizadora \textbf{WebPlotDigitizer} para realizar una calibración de ejes cartesianos y extraer con precisión milimétrica las coordenadas numéricas de casos semanales reales.

\section{Modelo simbólico}
El desarrollo del sistema integrado se compone de dos fases matemáticas sucesivas: un modelo predictivo basado en series temporales para estimar la demanda y un modelo de programación lineal entera mixta para optimizar la asignación de recursos.

\subsection{Modelo de Pronósticos (Holt-Winters Multiplicativo)}
La fase predictiva opera sobre un conjunto estructurado de índices y parámetros, donde $t$ denota la semana epidemiológica ($t = 1, 2, \dots, n$), $i$ representa la patología ($i = 1$ para Dengue, $i = 2$ para Bronquiolitis y $i = 3$ para Influenza), $s$ es la longitud del ciclo estacional anual ($s = 52$ semanas) y $m$ es el horizonte temporal de pronóstico.

El modelo de Suavización Exponencial Triple (Holt-Winters Multiplicativo) descompone la serie temporal real ($y_t$) actualizando semanalmente sus componentes a través de constantes de suavizamiento optimizadas en el intervalo $[0,1]$: de nivel ($\alpha$), de tendencia ($\beta$) y de estacionalidad ($\gamma$).

La ecuación de actualización del nivel base $L_t$ se formula como:
\begin{equation}
    L_t = \alpha \frac{y_t}{S_{t-s}} + (1 - \alpha) (L_{t-1} + T_{t-1})
\end{equation}
donde $L_t$ representa el nivel base, $y_t$ el valor real observado y $S_{t-s}$ el factor estacional del año previo. La tendencia secular $T_t$ se actualiza como:
\begin{equation}
    T_t = \beta (L_t - L_{t-1}) + (1 - \beta) T_{t-1}
\end{equation}
El factor de estacionalidad multiplicativo $S_t$ se calcula mediante:
\begin{equation}
    S_t = \gamma \frac{y_t}{L_t} + (1 - \gamma) S_{t-s}
\end{equation}
Finalmente, la proyección de la demanda provincial bruta para un horizonte futuro $m$ se formula como:
\begin{equation}
    F_{i,t+m} = (L_{t,i} + m T_{t,i}) S_{t-s+m,i}
\end{equation}
Para cuantificar la desviación y precisión de las proyecciones frente a los valores reales, se computan el Error Absoluto ($EA_t$) y la Desviación Absoluta Media ($MAD$):
\begin{equation}
    EA_t = |y_t - F_t|
\end{equation}
\begin{equation}
    MAD = \frac{1}{n} \sum_{t=1}^{n} |y_t - F_t|
\end{equation}

\subsection{Conversión de la Demanda y Segmentación Asistencial}
Para adaptar la demanda provincial bruta proyectada $F_{i,t+m}$ a la escala de planificación operativa del efector hospitalario nodal (Hospital Centro de Salud Zenón J. Santillán) en el modelo de optimización, se formaliza la conversión mediante la denominada ``Tabla de Uso'', fundamentada en dos coeficientes del sistema sanitario:

\begin{itemize}
    \item \textbf{Factor de absorción del subsistema público ($\theta_{\text{público}} = {,}60$):} Aunque la tasa de población sin cobertura formal de salud en Tucumán es del 42\% (Censo 2022 \cite{ref12}), durante contingencias epidémicas agudas (brotes de Dengue y picos de IRAB) el subsistema público (SIPROSA) absorbe aproximadamente el 60\% ($\theta_{\text{público}} = {,}60$) de la demanda directa efectiva. Esto responde a la saturación del sector privado, la falta de áreas de aislamiento febril o UTI pediátricas/adultas en clínicas privadas, y a la centralización de los protocolos de contingencia en la red pública.
    \item \textbf{Red de efectores de tercer nivel ($H_{\text{nodales}} = 6$):} La atención asistencial de contingencia de mediana y alta complejidad en la provincia se concentra en 6 hospitales públicos de Nivel III (Hospital Centro de Salud, Hospital Padilla, Hospital Avellaneda, Hospital del Niño Jesús, Hospital Eva Perón y Hospital de Concepción). La carga del sector público se distribuye equitativamente entre estos 6 efectores nodales.
\end{itemize}

La demanda hospitalaria semanal ajustada $D_i$ que ingresa como parámetro determinístico a las restricciones de flujo en la semana objetivo se calcula formalmente como:
\begin{equation}
    D_i = \left\lceil \frac{F_{i, t+m} \cdot \theta_{\text{público}}}{H_{\text{nodales}}} \right\rceil
\end{equation}
donde $\lceil \cdot \rceil$ representa la función techo para asegurar un valor entero y conservador. Para la Semana Epidemiológica 19, la demanda proyectada ajustada total asciende a 385 pacientes: $D_1 = 253$ (Dengue), $D_2 = 68$ (Bronquiolitis) y $D_3 = 64$ (Influenza).

\subsubsection{Segmentación por Modalidad Asistencial (Triaje Clínico)}
A fin de reflejar la realidad asistencial y evitar tratar a todos los pacientes como una categoría homogénea, la demanda total $D_i$ se divide en tres modalidades de atención: Cuidados Intensivos ($D_{i,\text{UTI}}$), Hospitalización General ($D_{i,\text{Gen}}$) y Atención Ambulatoria ($D_{i,\text{Amb}}$).

Se aplican las tasas de internación e ingreso a cuidados críticos por patología:
\begin{itemize}
    \item \textbf{Dengue ($i=1$):} 6{,}0\% de hospitalización base (0{,}5\% UTI, 5{,}5\% general, 94{,}0\% ambulatorio).
    \item \textbf{Bronquiolitis ($i=2$):} 12{,}5\% de hospitalización base (4{,}0\% UTI, 8{,}5\% general, 87{,}5\% ambulatorio).
    \item \textbf{Influenza ($i=3$):} 10{,}0\% de hospitalización base (5{,}0\% UTI, 5{,}0\% general, 90{,}0\% ambulatorio).
\end{itemize}

\textbf{Criterio de Redondeo Conservador:} Los pacientes de mayor gravedad ($D_{i,\text{UTI}}$ y $D_{i,\text{Gen}}$) se redondean hacia arriba mediante la función techo ($\lceil \cdot \rceil$) para evitar la subestimación de necesidades asistenciales críticas. La cantidad ambulatoria se obtiene por diferencia: $D_{i,\text{Amb}} = D_i - D_{i,\text{UTI}} - D_{i,\text{Gen}}$, garantizando que las tres modalidades sumen exactamente la demanda total $D_i$.

\begin{table}[htbp]
    \caption{Segmentación de la demanda hospitalaria por patología y modalidad (SE 19)}
    \label{tab:segmentacion_demanda}
    \centering
    \setlength{\tabcolsep}{3pt}
    \begin{tabular}{|>{\raggedright\arraybackslash}m{2.2cm}|>{\centering\arraybackslash}m{1.4cm}|>{\centering\arraybackslash}m{1.2cm}|>{\centering\arraybackslash}m{1.2cm}|>{\centering\arraybackslash}m{1.1cm}|}
        \hline
        \textbf{Patología ($i$)} & \textbf{Demanda Total ($D_i$)} & \textbf{UTI ($D_{i,\text{UTI}}$)} & \textbf{General ($D_{i,\text{Gen}}$)} & \textbf{Amb. ($D_{i,\text{Amb}}$)} \\
        \hline
        1. Dengue                & 253                            & 2                                 & 14                                    & 237                                \\
        \hline
        2. Bronquiolitis         & 68                             & 3                                 & 6                                     & 59                                 \\
        \hline
        3. Influenza             & 64                             & 4                                 & 3                                     & 57                                 \\
        \hline
        \textbf{Total}           & \textbf{385}                   & \textbf{9}                        & \textbf{23}                           & \textbf{353}                       \\
        \hline
    \end{tabular}
\end{table}

\subsection{Modelo de Optimización (Programación por Metas Lexicográfica por Subetapas)}
El modelo resuelve la asignación mediante una Programación por Metas Lexicográfica con prioridades absolutas ($P_1 \gg P_2$) para proteger la cobertura médica antes de evaluar la conveniencia financiera \cite{refGP}.

\subsubsection{Variables de Decisión}
\begin{itemize}
    \item $X_{i,m} \in \mathbb{Z}_0^+$: Pacientes atendidos de la patología $i$ en la modalidad $m \in \{\text{UTI}, \text{Gen}, \text{Amb}\}$ bajo el paquete completo exigido por guía técnica.
    \item $F_{i,m} \in \mathbb{Z}_0^+$: Pacientes no atendidos bajo el estándar completo de la guía técnica en la patología $i$ y modalidad $m$. No implica abandono del paciente, sino la activación de protocolos de contingencia (atención de contención sintomática en guardia, reprogramación de turnos y derivación a la red secundaria).
    \item $V_{r,i}$: Cantidad del recurso $r$ asignada a la patología $i$. Dominio discreto $V_{r,i} \in \mathbb{Z}_0^+$ para recursos indivisibles (camas en días-cama, kits, $r \in \{1,2,4,5\}$) y dominio continuo $V_{r,i} \in \mathbb{R}_0^+$ para recursos basados en tiempo (horas de oxígeno, médico, enfermería, $r \in \{3,6,7\}$).
    \item $d_2^-, d_2^+ \ge 0$: Variables de desvío de la meta financiera (presupuesto no utilizado y sobrecosto).
\end{itemize}

\subsubsection{Restricciones Estructurales (Duras)}
\begin{itemize}
    \item \textbf{Garantía de Protocolo Médico por Modalidad:} Un paciente solo se considera atendido si recibe el paquete completo de recursos requerido por protocolo:
          \begin{equation}
              V_{r,i} \ge \sum_{m \in \{\text{UTI}, \text{Gen}, \text{Amb}\}} T_{i,r,m} X_{i,m} \quad \forall r, i
          \end{equation}
    \item \textbf{Capacidad Neta Disponible:} Asignación sujeta a la capacidad nominal ($K_{\text{max},r}$) menos la carga base histórica ineludible ($B_r$):
          \begin{equation}
              \sum_{i=1}^{3} V_{r,i} \le K_{\text{neta},r} = K_{\text{max},r} - B_r \quad \forall r
          \end{equation}
          donde las camas generales y UTI se expresan en días-cama semanales ($K_{\text{max},1} = \text{camas físicas} \times 7 = 592{,}2$ días netos; $K_{\text{neta},2} = 79{,}8$ días netos de UTI).
    \item \textbf{Balance Epidemiológico de Flujo por Modalidad:}
          \begin{equation}
              X_{i,m} + F_{i,m} = D_{i,m} \quad \forall i, m
          \end{equation}
\end{itemize}

\subsubsection{Estructura Lexicográfica en Subetapas}
Para asegurar que los casos de mayor gravedad se atiendan antes de destinar recursos a menor complejidad, la Etapa 1 médica se subdivide en tres subetapas secuenciales, donde cada subetapa hereda una restricción de no retroceso:

\begin{enumerate}
    \item \textbf{Subetapa 1A (Cuidados Intensivos - UTI):} Minimiza el déficit de atención en pacientes críticos ponderado por gravedad clínica ($\mu_i$):
          \begin{equation}
              \text{Minimizar } d_{\text{UTI}}^- = \sum_{i=1}^{3} \mu_i F_{i,\text{UTI}}
          \end{equation}
          donde $\mu_1 = 1$ (Dengue), $\mu_2 = 2$ (Bronquiolitis) y $\mu_3 = 3$ (Influenza). El valor óptimo alcanzado es $d_{\text{UTI}}^{-*} = 0$.
    \item \textbf{Subetapa 1B (Hospitalización General):} Minimiza el déficit de internación general manteniendo la cobertura crítica inalterada:
          \begin{equation}
              \begin{aligned}
                  \text{Minimizar } & d_{\text{GEN}}^- = \sum_{i=1}^{3} \mu_i F_{i,\text{Gen}} \\
                  \text{s.a. }      & d_{\text{UTI}}^- \le d_{\text{UTI}}^{-*} = 0
              \end{aligned}
          \end{equation}
          El valor óptimo alcanzado es $d_{\text{GEN}}^{-*} = 0$.
    \item \textbf{Subetapa 1C (Atención Ambulatoria):} Maximiza la atención ambulatoria ponderada conservando en cero los déficits UTI y general:
          \begin{equation}
              \begin{aligned}
                  \text{Minimizar } & d_{\text{AMB}}^- = \sum_{i=1}^{3} \mu_i F_{i,\text{Amb}} \\
                  \text{s.a. }      & d_{\text{UTI}}^- = 0, \quad d_{\text{GEN}}^- = 0
              \end{aligned}
          \end{equation}
    \item \textbf{Etapa 2 (Prioridad Financiera):} Minimiza el sobrecosto $d_2^+$ respecto al fondo de contingencia asignado ($\text{Presupuesto} = \$17.000.000$), sin degradar ningún logro médico previo:
          \begin{equation}
              \begin{aligned}
                  \text{Minimizar } & d_2^+                                                                                          \\
                  \text{s.a. }      & d_{\text{UTI}}^- = 0, \quad d_{\text{GEN}}^- = 0, \quad d_{\text{AMB}}^- = d_{\text{AMB}}^{-*}
              \end{aligned}
          \end{equation}
          con la ecuación de meta financiera:
          \begin{equation}
              \sum_{r=1}^{7} \sum_{i=1}^{3} C_r V_{r,i} + d_2^- - d_2^+ = \text{Presupuesto}
          \end{equation}
\end{enumerate}




\section{Resolución}

\subsection{Fase Predictiva: Resultados y Ajuste del Pronóstico}
Para hallar la combinación óptima de las constantes de suavizamiento ($\alpha, \beta, \gamma$), se empleó el Solver de Microsoft Excel \cite{refExcel}, configurando como celda objetivo la minimización de la Desviación Absoluta Media ($MAD$) sobre los datos históricos.

El modelo de Holt-Winters Multiplicativo logró capturar las dinámicas específicas de cada patología:
\begin{itemize}
    \item \textbf{Dengue:} Mostró un comportamiento altamente explosivo y concentrado en verano y otoño, con factores estacionales muy marcados que multiplican el nivel base, típico de un brote vectorial agudo.
    \item \textbf{Bronquiolitis:} Evidenció picos predecibles durante los meses de invierno con baja dispersión interanual, facilitando la planificación de camas pediátricas y kits de ventilación respiratoria.
    \item \textbf{Influenza:} Reveló una circulación viral sostenida a lo largo del año con picos invernales superpuestos, permitiendo una proyección estable útil para prever horas de personal de salud y consumo de oxígeno.
\end{itemize}

A continuación se presentan las comparaciones gráficas entre las series históricas observadas y las proyecciones obtenidas mediante el modelo predictivo calibrado para las patologías evaluadas (Fig. \ref{fig:pronostico_dengue}, Fig. \ref{fig:pronostico_bronquiolitis} y Fig. \ref{fig:pronostico_influenza}):

\begin{figure}[htbp]
    \centering
    \includegraphics[width=\linewidth]{docs/Grafico_Dengue.png}
    \caption{Comparación de casos reales vs. pronóstico de Dengue en Tucumán (2024).}
    \label{fig:pronostico_dengue}
\end{figure}

\begin{figure}[htbp]
    \centering
    \includegraphics[width=\linewidth]{docs/Grafico_Bronquiolitis.png}
    \caption{Comparación de casos reales vs. pronóstico de Bronquiolitis en Tucumán (2024-2025).}
    \label{fig:pronostico_bronquiolitis}
\end{figure}

\begin{figure}[htbp]
    \centering
    \includegraphics[width=\linewidth]{docs/Grafico_Influenza.png}
    \caption{Comparación de casos reales vs. pronóstico de Influenza en Tucumán (2024-2025).}
    \label{fig:pronostico_influenza}
\end{figure}

\FloatBarrier
\subsection{Interfaz de Integración: Tabla de Uso}
El modelo de optimización por metas lexicográfica es de carácter estático y monoperíodo, diseñado para evaluar la planificación asistencial en una ventana temporal semanal aislada. Para articular este enfoque con la naturaleza dinámica y continua del modelo predictivo de series temporales, se implementó la herramienta denominada ``Tabla de Uso''.

Esta interfaz actúa como un filtro operativo que permite desacoplar la proyección secuencial de contagios y aislar semanas epidemiológicas críticas individuales (como el pico máximo de demanda de la SE 19). La Tabla de Uso toma la demanda bruta proyectada $F_{i,t+m}$, aplica los coeficientes de absorción pública ($\theta_{\text{público}} = {,}60$) y la división entre los efectores nodales ($H_{\text{nodales}} = 6$), e ingresa la demanda ajustada ($D_i = 385$) y su segmentación por triaje clínico como parámetros determinísticos estáticos en las restricciones del algoritmo de optimización.

\subsection{Fase de Optimización: Resultados de la GP Lexicográfica}
La resolución del modelo en dos etapas y tres subetapas arrojó los siguientes resultados:

\subsubsection{Etapa 1: Cobertura por Subetapas Médicas}
\begin{itemize}
    \item \textbf{Subetapa 1A (UTI):} Se atienden los 2 pacientes críticos de dengue, 3 de bronquiolitis y 4 de influenza. Cobertura UTI del 100\% (9/9), $d_{\text{UTI}}^- = 0$. Utiliza 62 días-cama UTI de los 79,8 disponibles (quedan 17,8 días-cama holgados).
    \item \textbf{Subetapa 1B (General):} Se atienden los 14 pacientes internados de dengue, 6 de bronquiolitis y 3 de influenza. Cobertura General del 100\% (23/23), $d_{\text{GEN}}^- = 0$. Utiliza 104 días-cama de los 592,2 disponibles.
    \item \textbf{Subetapa 1C (Ambulatoria):} Con los recursos remanentes, Solver atiende 68 pacientes ambulatorios de dengue, 0 de bronquiolitis y 12 de influenza. Déficit ambulatorio ponderado final $d_{\text{AMB}}^- = 422$ ($169 \times 1 + 59 \times 2 + 45 \times 3$).
\end{itemize}

En total, se atienden 112 de los 385 pacientes (29,1\% de cobertura global), garantizando el 100\% de protección a los casos graves e internados.

\begin{table}[htbp]
    \caption{Resumen asistencial de pacientes atendidos y no atendidos por modalidad}
    \label{tab:resultados_atencion}
    \centering
    \setlength{\tabcolsep}{3pt}
    \begin{tabular}{|>{\raggedright\arraybackslash}m{1.8cm}|>{\centering\arraybackslash}m{1.1cm}|>{\centering\arraybackslash}m{1.1cm}|>{\centering\arraybackslash}m{1.0cm}|>{\centering\arraybackslash}m{1.1cm}|>{\centering\arraybackslash}m{1.3cm}|}
        \hline
        \textbf{Patología} & \textbf{Amb.} & \textbf{General} & \textbf{UTI} & \textbf{Total Atend.} & \textbf{No Atend.} \\
        \hline
        Dengue             & 68            & 14               & 2            & 84                    & 169                \\
        \hline
        Bronquiolitis      & 0             & 6                & 3            & 9                     & 59                 \\
        \hline
        Influenza          & 12            & 3                & 4            & 19                    & 45                 \\
        \hline
        \textbf{Total}     & \textbf{80}   & \textbf{23}      & \textbf{9}   & \textbf{112}          & \textbf{273}       \\
        \hline
    \end{tabular}
\end{table}

\subsubsection{Uso de Recursos y Cuellos de Botella al Cierre de la Etapa 1}
El análisis de utilización de recursos demuestra que la infraestructura física pesada no constituye el factor limitante:

\begin{table}[htbp]
    \caption{Balance de utilización de recursos hospitalarios en la Etapa 1}
    \label{tab:uso_recursos}
    \centering
    \setlength{\tabcolsep}{3pt}
    \begin{tabular}{|>{\raggedright\arraybackslash}m{2.3cm}|>{\centering\arraybackslash}m{1.0cm}|>{\centering\arraybackslash}m{1.2cm}|>{\centering\arraybackslash}m{1.1cm}|>{\centering\arraybackslash}m{1.2cm}|}
        \hline
        \textbf{Recurso ($r$)} & \textbf{Uso} & \textbf{Cap. Neta} & \textbf{Holgura} & \textbf{\% Uso}  \\
        \hline
        Camas generales (días) & 104,0        & 592,2              & 488,2            & 17,6\%           \\
        \hline
        Camas UTI (días)       & 62,0         & 79,8               & 17,8             & 77,7\%           \\
        \hline
        Oxígeno/ARM (horas)    & 1.284,0      & 7.711,2            & 6.427,2          & 16,7\%           \\
        \hline
        Kits de hidratación    & 737,0        & 742,5              & 5,5              & 99,3\%           \\
        \hline
        Kits respiratorios     & 102,0        & 102,0              & 0,0              & \textbf{100,0\%} \\
        \hline
        Horas médicas          & 429,0        & 2.931,3            & 2.502,3          & 14,6\%           \\
        \hline
        Horas enfermería       & 1.510,0      & 4.167,6            & 2.657,6          & 36,2\%           \\
        \hline
    \end{tabular}
\end{table}

Los verdaderos cuellos de botella son los \textbf{Kits Respiratorios} (100,0\% de uso) y los \textbf{Kits de Hidratación} (99,3\% de uso). Estos insumos insustituibles agotan la capacidad asistencial ambulatoria a pesar de la holgura en camas y horas de personal.

\subsubsection{Etapa 2: Resultado Financiero y Sobrecosto}
La Etapa 2 minimiza el gasto operativo asegurando no retroceder en los resultados médicos alcanzados:
\begin{itemize}
    \item \textbf{Costo operativo mínimo requerido:} $\$40.321.805,93$.
    \item \textbf{Presupuesto contingencial disponible:} $\$17.000.000,00$.
    \item \textbf{Presupuesto no utilizado ($d_2^-$):} $\$0,00$.
    \item \textbf{Sobrecosto mínimo ($d_2^+$):} $\$23.321.805,93$.
\end{itemize}

El costo total equivale al 237,2\% del fondo asignado; el sobrecosto representa un 137,2\% adicional por encima del presupuesto. La mayor parte del costo operativo proviene de las Camas UTI (38,4\%) y Horas de Enfermería (21,9\%).

\section{Interpretación}

\subsection{Interpretación del Modelo de Pronósticos (Holt-Winters)}
El modelo de Holt-Winters Multiplicativo captura con precisión la dinámica epidemiológica de Tucumán. La anulación de la constante de tendencia ($\beta \to 0$) por parte de Solver confirma que las epidemias estacionales presentan un comportamiento oscilatorio cíclico anual sin tendencia secular de largo plazo.

\subsection{Interpretación del Modelo de Optimización (GP Lexicográfica por Subetapas)}
La estructura en subetapas resuelve de forma éticamente irreprochable el triaje:
\begin{enumerate}
    \item Protege al 100\% de las vidas en riesgo crítico (UTI) y de los casos graves con necesidad de hospitalización general.
    \item Evita la monetización de la salud al impedir que la Meta 2 altere la cobertura alcanzada en la Meta 1.
    \item Revela de manera objetiva que la brecha ambulatoria no se debe a restricciones de espacio ni de personal, sino a cuellos de botella logísticos en insumos descartables (kits).
    \item Identifica un dilema de política sanitaria en la asignación ambulatoria: Bronquiolitis recibe 0 atención ambulatoria en el óptimo matemático puro debido a que Influenza posee un ponderador de severidad mayor ($\mu_3 = 3$ vs. $\mu_2 = 2$) y consume menos kits respiratorios por paciente. Este hallazgo demuestra la necesidad de que los gestores sanitarios complementen la optimización pura con políticas explícitas de equidad.
    \item Clarifica el significado sanitario de la variable de déficit $F_{i,m}$ (pacientes no atendidos bajo el estándar completo del efector): un valor $F_{i,m} > 0$ no implica de ningún modo el abandono ni que se deje a la deriva al paciente. En el contexto de la gestión sanitaria pública (SIPROSA), indica únicamente que no es posible dispensar el paquete de atención integral con la rigurosidad técnica estándar fijada por las guías clínicas. La resolución asistencial para esta demanda no absorbida en el estándar primario se encauza mediante un esquema tripartito: a) atención sintomática de contención en guardia; b) reprogramación de turnos no urgentes; y c) derivación institucional coordinada hacia efectores de la red secundaria de salud.
\end{enumerate}

\section{Conclusiones}

El presente trabajo ha permitido establecer un marco analítico integral de soporte a las decisiones hospitalarias ante crisis epidemiológicas, articulando pronósticos de series temporales y optimización multiobjetivo. A partir del desarrollo e implementación de la metodología propuesta, se arriban a las siguientes conclusiones principales:

\begin{itemize}
    \item \textbf{Modelación Predictiva de Brotes Epidémicos:} La Suavización Exponencial Triple de Holt-Winters con estacionalidad multiplicativa demostró ser una herramienta altamente efectiva para caracterizar patrones no lineales, logrando capturar tanto los picos epidémicos del dengue como el comportamiento cíclico invernal de las infecciones respiratorias agudas bajas (bronquiolitis e influenza). Asimismo, la convergencia a cero de la constante de tendencia ($\beta \to 0$) valida la hipótesis epidemiológica de que los brotes estacionales responden a dinámicas oscilatorias anuales reguladas por factores estacionales y niveles base, sin tendencias seculares persistentes de largo plazo.
    \item \textbf{Protección Asistencial y Rigor Ético-Médico:} La estructuración del modelo de Programación por Metas Lexicográfica en subetapas secuenciales resolvió de manera robusta el triaje clínico. Se garantizó la cobertura integral (100\%) de los pacientes en estado crítico (UTI) y de aquellos que requerían hospitalización general, demostrando que la prioridad asistencial médica puede preservarse de forma innegociable sin supeditar el valor de la vida humana a restricciones presupuestarias.
    \item \textbf{Interpretación Sanitaria de la Demanda No Atendida ($F_{i,m}$):} Se estableció que un registro de déficit $F_{i,m} > 0$ no constituye abandono ni desatención médica del paciente, sino la imposibilidad puntual de aplicar el protocolo asistencial bajo la máxima rigurosidad prescrita por las guías técnicas internacionales. La red pública de salud (SIPROSA) gestiona esta demanda mediante mecanismos de contención inicial en guardia, reprogramación oportuna de consultas ambulatorias y derivación asistida hacia efectores secundarios.
    \item \textbf{Identificación de Cuellos de Botella Operativos:} El análisis de capacidad evidenció que la restricción principal del sistema público de salud ante escenarios de contingencia no se encuentra en la infraestructura física rígida de camas (17,6\% de uso en sala general y 77,7\% en UTI) ni en la disponibilidad de personal sanitario (14,6\% de horas médicas y 36,2\% de horas de enfermería). Por el contrario, los verdaderos cuellos de botella críticos radican en el desabastecimiento de insumos descartables específicos, como Kits Respiratorios (100,0\% de utilización) y Kits de Hidratación (99,3\%), ofreciendo a los gestores de salud una directriz operativa clara para focalizar las políticas de compra y stock contingencial.
    \item \textbf{Cuantificación del Déficit Financiero:} La segunda etapa del modelo permitió determinar con precisión el costo operativo mínimo necesario para sostener la cobertura médica priorizada (\$40.321.805,93), identificando una brecha presupuestaria del 137,2\% sobre el fondo de contingencia asignado (\$17.000.000,00). Esta cuantificación cuantitativa ofrece una base objetiva para justificar reasignaciones presupuestarias ministeriales ante emergencias sanitarias.
    \item \textbf{Aporte Metodológico y Resguardo Ético:} El sistema desarrollado constituye un instrumento de gestión proactivo, transparente y orientativo para la red de salud pública, capaz de prever la asignación óptima de recursos sin reemplazar jamás el juicio clínico. Se enfatiza que el modelo matemático actúa como soporte cuantitativo de planificación, permaneciendo siempre supeditado al criterio médico profesional, la deontología sanitaria y los valores éticos y morales de la sociedad.
    \item \textbf{Trabajo Futuro y Extensión Dinámica Multiperíodo:} Como propuesta de trabajo futuro, se plantea extender la presente formulación estática monoperíodo hacia un modelo dinámico multiperíodo que incorpore ecuaciones de balance de inventario de insumos entre semanas epidemiológicas consecutivas, traslape de la duración real de estancia hospitalaria ($LOS$) y programación estocástica o por escenarios para cuantificar formalmente la incertidumbre de la demanda proyectada.
\end{itemize}

\vfill\newpage
\begin{thebibliography}{00}
    \bibitem{ref1} Ministerio de Salud Pública de Tucumán. (2024). Sala de situación epidemiológica. Dirección de Epidemiología.
    \bibitem{ref2} Mahmud, N., Ali, N. A., Pazil, N. S. M., \& Jamaluddin, S. H. (2022). Predicting dengue outbreak in Selangor using Holt-Winters models. International Journal of Academic Research in Business and Social Sciences, 12(1).
    \bibitem{ref3} Mahmud, N., \& Pazil, N. S. M. (2021). Prediction of dengue outbreak: A comparison between ARIMA and Holt-Winters methods. Universiti Teknologi MARA.
    \bibitem{ref4} Shah, D., Marathe, S. A., \& Mehrotra, M. (2023). Operations research applications in healthcare systems: Resource allocation optimization models. Journal of Applied Analytics and Forecasting Research.
    \bibitem{ref5} Shams Eddin, M., \& El Hajj, H. (2025). An optimization-based framework to dynamically schedule hospital beds in a pandemic. Preprints.
    \bibitem{ref6} Dirección de Investigación para la Salud. (2017). Agenda Nacional de Investigación en Salud Pública: Diagnóstico de Situación de Áreas de Vacancia en Investigación Operativa en Salud de Argentina. Ministerio de Salud de la Nación.
    \bibitem{ref7} Anderson, D. R., Sweeney, D. J., Williams, T. A., Camm, J. D., \& Martin, K. (2011). Métodos cuantitativos para los negocios (11.ª ed.). CENGAGE Learning.
    \bibitem{ref11} Winston, W. L. (2005). Investigación de Operaciones: Aplicaciones y algoritmos (4.ª ed.). THOMSON.
    \bibitem{ref9} Ministerio de Salud de la Nación Argentina. (2023). Estrategia de Gestión Integrada para la Prevención y Control del Dengue. Dirección Nacional de Epidemiología y Análisis de Situación de Salud.
    \bibitem{ref10} Sistema Provincial de Salud de Tucumán (SIPROSA). (2024). Plan de contingencia y lineamientos operativos para la readecuación de recursos físicos ante brotes epidémicos. Ministerio de Salud Pública, Gobierno de Tucumán.
    \bibitem{ref8} Bonomo, F., Durán, G., \& Marenco, J. (2020). Modelos de simulación y optimización para la gestión de recursos hospitalarios ante variaciones de demanda. Instituto de Cálculo, Universidad de Buenos Aires - CONICET.
    \bibitem{ref12} Instituto Nacional de Estadística y Censos [INDEC]. (2023). Censo Nacional de Población, Hogares y Viviendas 2022: Resultados definitivos. Indicadores demográficos por provincia.
    \bibitem{ref13} Ministerio de Salud de la Nación. (2024). Boletín Epidemiológico Nacional N° 702 - SE 17.
    \bibitem{ref14} Organización Mundial de la Salud [OMS]. (2023). Dengue y dengue grave: Datos y cifras.
    \bibitem{ref15} Organización Panamericana de la Salud [OPS]. (2022). Actualización Epidemiológica: Virus Respiratorio Sincitial (VSR) e Influenza en la Región de las Américas.
    \bibitem{refExcel} Veliz, I. M., Bellor, M. E., Diaz, D. S., Jimenez, F. E., \& Zato, C. A. (2026). Planilla de soporte y calibración del modelo de pronósticos: \emph{Optimización de la gestión de recursos críticos hospitalarios}. Disponible en: \url{https://bit.ly/pronosticos_invop}
    \bibitem{refGP} Veliz, I. M., Bellor, M. E., Diaz, D. S., Jimenez, F. E., \& Zato, C. A. (2026). Modelo de Asignación de Recursos Hospitalarios — GP Lexicográfica. \emph{Optimización de la gestión de recursos críticos hospitalarios}. Disponible en: \url{http://bit.ly/modelo_gp_invop}
    \bibitem{refMSPT2} Ministerio de Salud Pública de Tucumán. \emph{Hospital Centro de Salud Zenón J. Santillán: Características y servicios}. Gobierno de Tucumán. Disponible en: \url{https://msptucuman.gov.ar/hospitales/hospital-centro-de-salud-zenon-j-santillan/caracteristicas/}
    \bibitem{refLey5908} Poder Ejecutivo de la Provincia de Tucumán. \emph{Ley Nº 5908: Carrera Sanitaria Provincial y modificatorias}. Ministerio de Salud Pública de Tucumán. Disponible en: \url{https://msptucuman.gov.ar/wordpress/wp-content/uploads/2023/12/Ley-5908-Carrera-Sanitaria-original-Fue-modificada-por-Ley-8757.pdf}
    \bibitem{refATE} Asociación Trabajadores del Estado (ATE). \emph{Escalas salariales y convenios colectivos del sector salud (Decreto Nº 1133/09)}. Convenios Colectivos ATE. Disponible en: \url{https://convenios.ate.org.ar/wp-content/uploads/2026/06/Salud-1133-jun-ago2026.pdf}

    \bibitem{refOPS2} Organización Panamericana de la Salud (OPS). \emph{Documentos técnicos y directrices regionales en salud pública}. OPS/Iris. Disponible en: \url{https://iris.paho.org/items/0832bcad-f6b2-4c52-bb28-eac8f587fb6e}
    \bibitem{refOMS2} Organización Mundial de la Salud (OMS). \emph{Publicaciones científicas y técnicas institucionales}. OMS/Iris. Disponible en: \url{https://iris.who.int/items/1c1f7fd3-44cd-4029-9a95-4470427834bf}
\end{thebibliography}

\end{document}